% GENERATED FILE. Edit canonical lessons, then run npm run generate:books.
% canonical-source-hash: fnv1a64:877c4d95ca69d157
% canonical-lessons: GU-C542-samadhan, GU-C542-vataghat, GU-C542-madhyasthi, GU-C542-bhalaman, GU-C542-vinanti

\chapter{Meeting Halfway}
\label{ch:gu-meeting-halfway}
\hlchaptermodality{542}

\begin{chapteropening}
\textbf{By the end of this chapter:} \emph{I can talk about talks, compromise and asking someone for help.}

Complete the last lesson of chapter 542: I can talk about talks, compromise and asking someone for help.
\end{chapteropening}

\section[\texorpdfstring{\gu{સમાધાન} (samādhān)}{સમાધાન (samādhān)}]{\gu{સમાધાન} (samādhān) — a compromise}

\label{lesson:GU-C542-samadhan}



\begin{quote}
Before the new one: say the Gujarati for a loophole, then the Gujarati for a proposal.
\end{quote}

\subsection*{You'll want to know: \gu{સમાધાન}}
\textbf{\gu{સમાધાન}} — \emph{samādhān} — \textquotedblleft{}a compromise\textquotedblright{}.

\begin{quote}
\gu{જો} \gu{તમે} \gu{બંને} \gu{બેસીને} \gu{વાત} \gu{કરશો} \gu{તો} \gu{સમાધાન} \gu{થશે}.
\end{quote}

\begin{quote}
\emph{jo tame banne besīne vāt karśo to samādhān thaśe.}
\end{quote}

\begin{quote}
If you two sit down and talk, you'll reach a compromise.
\end{quote}

A neuter noun.

\begin{cousinweb}
\textbf{\gu{સમાધાન}} (\emph{samādhān}) is a tatsama, taken from Sanskrit \emph{samādhāna}, \textquotedblleft{}putting together\textquotedblright{} (Monier-Williams); in Gujarati it is the settling of a quarrel.
\end{cousinweb}

\begin{grammarlens}[title={What you've built}]
One more word for saying what happens if.
\end{grammarlens}

\subsection*{Your turn}
\begin{itemize}
\raggedright
  \item \emph{Say it:} \emph{samādhān}
  \item \emph{Say it:} \emph{samādhān}, once more
  \item \emph{Say it:} \emph{jo tame banne besīne vāt karśo to samādhān thaśe.} — If you two sit down and talk, you'll reach a compromise.
  \item \emph{Recall:} say the Gujarati for a loophole, then the Gujarati for a proposal, then say \emph{samādhān} again
\end{itemize}

\begin{tcolorbox}[breakable,colback=teal!4,colframe=teal!35!black,title={Before you move on}]
What does \gu{સમાધાન} mean? (\textquotedblleft{}A compromise\textquotedblright{}.) Say it once more.
\end{tcolorbox}
\section[\texorpdfstring{\gu{વાટાઘાટ} (vāṭāghāṭ)}{વાટાઘાટ (vāṭāghāṭ)}]{\gu{વાટાઘાટ} (vāṭāghāṭ) — negotiations}

\label{lesson:GU-C542-vataghat}



\begin{quote}
Before the new one: say the Gujarati for a proposal, then the Gujarati for a compromise.
\end{quote}

\subsection*{You'll want to know: \gu{વાટાઘાટ}}
\textbf{\gu{વાટાઘાટ}} — \emph{vāṭāghāṭ} — \textquotedblleft{}negotiations\textquotedblright{}.

\begin{quote}
\gu{જો} \gu{વાટાઘાટ} \gu{સારી} \gu{ચાલશે} \gu{તો} \gu{કાલે} \gu{કારખાનું} \gu{ફરી} \gu{ખૂલશે}.
\end{quote}

\begin{quote}
\emph{jo vāṭāghāṭ sārī cālśe to kāle kārkhānũ pharī khūlśe.}
\end{quote}

\begin{quote}
If the talks go well, the factory will open again tomorrow.
\end{quote}

A feminine noun.

\begin{cousinweb}
\textbf{\gu{વાટાઘાટ}} (\emph{vāṭāghāṭ}) is a jingling pair of words for talk back and forth between two sides. No dictionary we use records an older form, so its history is not settled.
\end{cousinweb}

\begin{grammarlens}[title={What you've built}]
One more word for saying what happens if.
\end{grammarlens}

\subsection*{Your turn}
\begin{itemize}
\raggedright
  \item \emph{Say it:} \emph{vāṭāghāṭ}
  \item \emph{Say it:} \emph{vāṭāghāṭ}, once more
  \item \emph{Say it:} \emph{jo vāṭāghāṭ sārī cālśe to kāle kārkhānũ pharī khūlśe.} — If the talks go well, the factory will open again tomorrow.
  \item \emph{Recall:} say the Gujarati for a proposal, then the Gujarati for a compromise, then say \emph{vāṭāghāṭ} again
\end{itemize}

\begin{tcolorbox}[breakable,colback=teal!4,colframe=teal!35!black,title={Before you move on}]
What does \gu{વાટાઘાટ} mean? (\textquotedblleft{}Negotiations\textquotedblright{}.) Say it once more.
\end{tcolorbox}
\section[\texorpdfstring{\gu{મધ્યસ્થી} (madhyasthī)}{મધ્યસ્થી (madhyasthī)}]{\gu{મધ્યસ્થી} (madhyasthī) — a mediator}

\label{lesson:GU-C542-madhyasthi}



\begin{quote}
Before the new one: say the Gujarati for a compromise, then the Gujarati for negotiations.
\end{quote}

\subsection*{You'll want to know: \gu{મધ્યસ્થી}}
\textbf{\gu{મધ્યસ્થી}} — \emph{madhyasthī} — \textquotedblleft{}a mediator\textquotedblright{}.

\begin{quote}
\gu{જો} \gu{કોઈ} \gu{મધ્યસ્થી} \gu{મદદ} \gu{કરશે} \gu{તો} \gu{બંને} \gu{ભાઈ} \gu{ફરી} \gu{મળશે}.
\end{quote}

\begin{quote}
\emph{jo koī madhyasthī madad karśe to banne bhāī pharī maḷśe.}
\end{quote}

\begin{quote}
If a mediator helps, the two brothers will meet again.
\end{quote}

A masculine or feminine noun, by the person meant.

\begin{cousinweb}
\textbf{\gu{મધ્યસ્થી}} (\emph{madhyasthī}) is built on the tatsama Sanskrit \emph{madhyastha}, \textquotedblleft{}standing in the middle\textquotedblright{} (Monier-Williams): the one who stands between two sides.
\end{cousinweb}

\begin{grammarlens}[title={What you've built}]
One more word for saying what happens if.
\end{grammarlens}

\subsection*{Your turn}
\begin{itemize}
\raggedright
  \item \emph{Say it:} \emph{madhyasthī}
  \item \emph{Say it:} \emph{madhyasthī}, once more
  \item \emph{Say it:} \emph{jo koī madhyasthī madad karśe to banne bhāī pharī maḷśe.} — If a mediator helps, the two brothers will meet again.
  \item \emph{Recall:} say the Gujarati for a compromise, then the Gujarati for negotiations, then say \emph{madhyasthī} again
\end{itemize}

\begin{tcolorbox}[breakable,colback=teal!4,colframe=teal!35!black,title={Before you move on}]
What does \gu{મધ્યસ્થી} mean? (\textquotedblleft{}A mediator\textquotedblright{}.) Say it once more.
\end{tcolorbox}
\section[\texorpdfstring{\gu{ભલામણ} (bhalāmaṇ)}{ભલામણ (bhalāmaṇ)}]{\gu{ભલામણ} (bhalāmaṇ) — a recommendation}

\label{lesson:GU-C542-bhalaman}



\begin{quote}
Before the new one: say the Gujarati for negotiations, then the Gujarati for a mediator.
\end{quote}

\subsection*{You'll want to know: \gu{ભલામણ}}
\textbf{\gu{ભલામણ}} — \emph{bhalāmaṇ} — \textquotedblleft{}a recommendation\textquotedblright{}.

\begin{quote}
\gu{જો} \gu{સાહેબ} \gu{ભલામણ} \gu{કરશે} \gu{તો} \gu{મને} \gu{નોકરી} \gu{મળશે}.
\end{quote}

\begin{quote}
\emph{jo sāheb bhalāmaṇ karśe to mane nokrī maḷśe.}
\end{quote}

\begin{quote}
If the boss recommends me, I'll get the job.
\end{quote}

A feminine noun.

\begin{cousinweb}
\textbf{\gu{ભલામણ}} (\emph{bhalāmaṇ}) is a good word put in for someone. No dictionary we use records an older form, so its history is not settled.
\end{cousinweb}

\begin{grammarlens}[title={What you've built}]
One more word for saying what happens if.
\end{grammarlens}

\subsection*{Your turn}
\begin{itemize}
\raggedright
  \item \emph{Say it:} \emph{bhalāmaṇ}
  \item \emph{Say it:} \emph{bhalāmaṇ}, once more
  \item \emph{Say it:} \emph{jo sāheb bhalāmaṇ karśe to mane nokrī maḷśe.} — If the boss recommends me, I'll get the job.
  \item \emph{Recall:} say the Gujarati for negotiations, then the Gujarati for a mediator, then say \emph{bhalāmaṇ} again
\end{itemize}

\begin{tcolorbox}[breakable,colback=teal!4,colframe=teal!35!black,title={Before you move on}]
What does \gu{ભલામણ} mean? (\textquotedblleft{}A recommendation\textquotedblright{}.) Say it once more.
\end{tcolorbox}
\section[\texorpdfstring{\gu{વિનંતી} (vinantī)}{વિનંતી (vinantī)}]{\gu{વિનંતી} (vinantī) — a polite request}

\label{lesson:GU-C542-vinanti}



\begin{quote}
Before the new one: say the Gujarati for a mediator, then the Gujarati for a recommendation.
\end{quote}

\subsection*{You'll want to know: \gu{વિનંતી}}
\textbf{\gu{વિનંતી}} — \emph{vinantī} — \textquotedblleft{}a polite request\textquotedblright{}.

\begin{quote}
\gu{જો} \gu{તમે} \gu{વિનંતી} \gu{કરશો} \gu{તો} \gu{તે} \gu{ચોક્કસ} \gu{મદદ} \gu{કરશે}.
\end{quote}

\begin{quote}
\emph{jo tame vinantī karśo to te cokkas madad karśe.}
\end{quote}

\begin{quote}
If you ask him politely, he will surely help.
\end{quote}

A feminine noun.

\begin{cousinweb}
\textbf{\gu{વિનંતી}} (\emph{vinantī}) goes back to Sanskrit \emph{vijñapti}, \textquotedblleft{}a report; a request to a superior\textquotedblright{}, and CDIAL lists the Gujarati word, \textquotedblleft{}a request\textquotedblright{} (CDIAL 11706).
\end{cousinweb}

\begin{grammarlens}[title={What you've built}]
One more word for saying what happens if.
\end{grammarlens}

\subsection*{Your turn}
\begin{itemize}
\raggedright
  \item \emph{Say it:} \emph{vinantī}
  \item \emph{Say it:} \emph{vinantī}, once more
  \item \emph{Say it:} \emph{jo tame vinantī karśo to te cokkas madad karśe.} — If you ask him politely, he will surely help.
  \item \emph{Recall:} say the Gujarati for a mediator, then the Gujarati for a recommendation, then say \emph{vinantī} again
\end{itemize}

\begin{tcolorbox}[breakable,colback=teal!4,colframe=teal!35!black,title={Before you move on}]
What does \gu{વિનંતી} mean? (\textquotedblleft{}A polite request\textquotedblright{}.) Say it once more.
\end{tcolorbox}
